\documentclass[twocolumn,prl,amsmath,amssymb,superscriptaddress]{revtex4-2}

\usepackage{graphicx}
\usepackage{hyperref}
\usepackage{amsmath}
\usepackage{amssymb}
\usepackage{booktabs}
\usepackage{algorithm}
\usepackage{algpseudocode}
\usepackage{xcolor}
\usepackage{tikz}

\newcommand{\zxmotifs}{\textsc{zx-motifs}}
\newcommand{\Ket}[1]{\left|#1\right\rangle}

\begin{document}

\title{ZX-Motifs: Structural Fingerprinting and Data-Driven Discovery of Quantum Algorithms via ZX-Calculus Subgraph Analysis}

\author{Spencer Churchill}
\affiliation{University of California, Santa Cruz}

\date{\today}

\begin{abstract}
We introduce \zxmotifs{}, an open-source framework that bridges quantum circuit analysis and graph-theoretic pattern mining through the ZX-calculus.
The system converts a corpus of 78 quantum algorithms spanning 19 algorithmic families into ZX-calculus diagrams at six levels of diagrammatic simplification, extracts semantically labeled graph representations, and discovers recurring structural motifs via subgraph isomorphism.
A three-pronged motif discovery strategy---combining domain-informed templates, bottom-up subgraph enumeration, and neighborhood extraction---yields 697 unique motifs after Weisfeiler--Leman deduplication.
Fingerprint vectors constructed from motif occurrence counts enable hierarchical clustering that recovers known algorithmic family structure and reveals cross-family structural kinships.
We identify six universal motifs present in over 80\% of algorithm families and track motif survival through simplification levels, finding 10 \emph{irreducible} motifs that persist under full ZX reduction.
Composing pairs of irreducible motifs yields novel circuit ansätze; one such ansatz, \texttt{irr\_pair11}, statistically significantly outperforms hardware-efficient and linear-chain baselines across six VQE benchmark configurations (p < 0.001, Cohen's $d$ up to 450).
Ablation studies confirm that each structural component---star-hub entanglement, non-Clifford phase gadgets, and chain-tail connectivity---contributes to the ansatz's expressibility.
Our results demonstrate that ZX-calculus motif analysis provides a principled, automated approach to quantum algorithm taxonomy and structure-guided circuit design.
\end{abstract}

\maketitle

% ============================================================================
\section{Introduction}
\label{sec:introduction}

The rapidly growing landscape of quantum algorithms presents a classification challenge: algorithms developed for disparate applications---from optimization and error correction to quantum simulation and machine learning---share underlying structural primitives that are obscured by their gate-level descriptions.
Understanding these shared structures is valuable both for algorithm taxonomy and for the design of new circuits.

The ZX-calculus~\cite{Coecke2011,vandeWetering2020} provides a natural framework for this analysis.
By representing quantum circuits as string diagrams composed of Z-spiders, X-spiders, and Hadamard edges, the ZX-calculus abstracts away gate-level details while preserving the essential algebraic structure.
Importantly, ZX diagrams admit a hierarchy of simplification rules---from spider fusion through Clifford simplification to full reduction---that progressively reveal deeper structural properties.

In this work, we develop \zxmotifs{}, a complete pipeline for converting quantum circuits into ZX diagrams, discovering recurring structural patterns (motifs), and using these motifs for algorithm fingerprinting, classification, and data-driven circuit design.
Our contributions are:

\begin{enumerate}
    \item A comprehensive corpus of 78 quantum algorithms across 19 families, each convertible to ZX diagrams at six simplification levels (Section~\ref{sec:corpus}).
    \item A three-strategy motif discovery framework combining hand-crafted templates, bottom-up enumeration, and neighborhood extraction, yielding 697 unique motifs (Section~\ref{sec:discovery}).
    \item A fingerprinting and phylogenetic analysis that recovers known algorithm families and reveals unexpected structural relationships (Section~\ref{sec:fingerprinting}).
    \item A structure-guided ansatz discovery method that composes irreducible motifs into novel variational circuits, with one ansatz achieving statistically significant VQE improvements (Section~\ref{sec:ansatz_discovery}).
\end{enumerate}

% ============================================================================
\section{Background}
\label{sec:background}

\subsection{ZX-Calculus}

The ZX-calculus represents quantum processes as open graphs with two generator types: \emph{Z-spiders} (green nodes) and \emph{X-spiders} (red nodes), each carrying a phase $\alpha \in [0, 2\pi)$, connected by simple (identity) or Hadamard edges.
A Z-spider with phase $\alpha$ and $m$ inputs, $n$ outputs implements the linear map
\begin{equation}
    \Ket{0}^{\otimes n}\Bra{0}^{\otimes m} + e^{i\alpha}\Ket{1}^{\otimes n}\Bra{1}^{\otimes m},
\end{equation}
with X-spiders defined analogously in the $\Ket{\pm}$ basis.
The calculus is \emph{complete} for qubit quantum mechanics~\cite{Hadzihasanovic2018,Jeandel2018}, meaning any equation between quantum processes provable by matrix multiplication is also derivable from the ZX rewrite rules.

\subsection{Simplification Hierarchy}

PyZX~\cite{Kissinger2020} implements a hierarchy of simplification strategies that we exploit to reveal structure at multiple abstraction levels:

\begin{enumerate}
    \item \textbf{RAW}: Direct translation from QASM; no simplification.
    \item \textbf{SPIDER\_FUSED}: Adjacent same-color spiders fused; phases combined.
    \item \textbf{INTERIOR\_CLIFFORD}: Interior Clifford spiders eliminated via local complementation and pivoting.
    \item \textbf{CLIFFORD\_SIMP}: Full Clifford simplification.
    \item \textbf{FULL\_REDUCE}: Iterative application of all reduction rules.
    \item \textbf{TELEPORT\_REDUCE}: Phase teleportation followed by reduction.
\end{enumerate}

Each level subsumes the previous, progressively compressing the diagram while preserving the represented unitary (up to global phase and ancilla introduction at the highest levels).

\subsection{Subgraph Isomorphism}

Motif detection reduces to the subgraph isomorphism problem: given a small \emph{pattern} graph $P$ and a larger \emph{host} graph $H$, find all injective mappings $\phi: V(P) \to V(H)$ that preserve adjacency and all semantic labels.
We employ the VF2 algorithm~\cite{Cordella2004} as implemented in NetworkX, augmented with custom node and edge matching predicates.

% ============================================================================
\section{System Architecture}
\label{sec:architecture}

\zxmotifs{} implements a modular pipeline:
\begin{equation*}
    \text{Qiskit} \xrightarrow{\text{QASM2}} \text{PyZX} \xrightarrow{\text{feature}} \text{NetworkX} \xrightarrow{\text{match}} \text{Fingerprint}
\end{equation*}

\subsection{Conversion and Featurization}

Each Qiskit \texttt{QuantumCircuit} is exported to QASM~2.0 and parsed by PyZX into a ZX graph.
The graph is then simplified at all six levels, producing a \texttt{ZXSnapshot} that records the graph, vertex/edge counts, T-gate count, and a circuit-likeness heuristic.

The featurizer converts each PyZX graph into a labeled NetworkX graph with the following semantic attributes:

\textbf{Node attributes:}
\begin{itemize}
    \item \texttt{vertex\_type} $\in \{Z, X, H\_BOX, BOUNDARY\}$
    \item \texttt{phase\_class} $\in \{$zero, pauli, clifford, t\_like, arbitrary$\}$
    \item \texttt{degree} (integer)
\end{itemize}

\textbf{Edge attributes:}
\begin{itemize}
    \item \texttt{edge\_type} $\in \{SIMPLE, HADAMARD\}$
\end{itemize}

Phase coarsening maps exact phases to equivalence classes: $0 \mapsto$ zero, $\pi \mapsto$ pauli, $k\pi/2 \mapsto$ clifford, $k\pi/4 \mapsto$ t\_like, else $\mapsto$ arbitrary.
This abstraction enables parametric motif matching while preserving algebraically meaningful distinctions.

\subsection{Semantic Subgraph Matching}

The matcher performs VF2 subgraph isomorphism with semantic constraints.
Two nodes match if they share the same \texttt{vertex\_type} and compatible \texttt{phase\_class}; two edges match if they share the same \texttt{edge\_type}.

Critically, we introduce \emph{phase wildcards} for parametric motifs:
\begin{itemize}
    \item \texttt{any}: matches all phase classes
    \item \texttt{any\_nonzero}: matches pauli, clifford, t\_like, arbitrary
    \item \texttt{any\_nonclifford}: matches t\_like, arbitrary
\end{itemize}

A pre-filter (\texttt{can\_possibly\_match}) checks necessary conditions---type-count dominance, edge-count inequality, and degree-sequence feasibility---to prune the search space before invoking VF2.

% ============================================================================
\section{Algorithm Corpus}
\label{sec:corpus}

We assembled a corpus of 78 quantum algorithms across 19 families (Table~\ref{tab:families}).
Each algorithm is implemented as a parameterized Qiskit generator function that accepts a qubit count $n$ and returns a \texttt{QuantumCircuit}.
Algorithms span the major application domains of quantum computing:

\begin{table}[h]
\caption{Algorithm families in the \zxmotifs{} corpus. Each family contains parameterized circuit generators for multiple qubit counts.}
\label{tab:families}
\begin{ruledtabular}
\begin{tabular}{lrl}
\textbf{Family} & \textbf{Count} & \textbf{Representative} \\
\midrule
Variational & 10 & QAOA, VQE-UCCSD, ADAPT-VQE \\
Oracle & 9 & Grover, Bernstein--Vazirani \\
Error Correction & 9 & Steane, Surface Code \\
Simulation & 8 & Trotter--Ising, qDRIFT \\
Entanglement & 6 & GHZ, W-state, Cluster \\
Machine Learning & 6 & QSVM, QCNN, QGAN \\
Protocol & 4 & Teleportation, E91 \\
Transform & 4 & QFT, Phase Estimation \\
Arithmetic & 4 & Ripple-carry, QFT Adder \\
Distillation & 4 & BBPSSW, DEJMPS \\
Cryptography & 2 & BB84, E91 \\
Sampling & 2 & IQP, Random Circuit \\
Error Mitigation & 2 & ZNE, Pauli Twirling \\
Topological & 2 & Jones Polynomial, Toric Code \\
Metrology & 2 & GHZ Metrology, Fisher Info \\
Linear Algebra & 2 & HHL, VQLS \\
Diff.\ Equations & 1 & Poisson Solver \\
TDA & 1 & Betti Number \\
Communication & 1 & Fingerprinting \\
\end{tabular}
\end{ruledtabular}
\end{table}

Conversion to ZX diagrams at all six simplification levels yields 1,104 total graph instances (184 algorithm instances $\times$ 6 levels).
We observe substantial compression: from an average of 72.4 vertices at RAW level to 30.9 vertices at FULL\_REDUCE (57\% reduction), confirming that significant redundancy exists in gate-level descriptions.

% ============================================================================
\section{Motif Discovery}
\label{sec:discovery}

We employ three complementary strategies for motif discovery, each targeting different aspects of the structural landscape.

\subsection{Hand-Crafted Motifs}

Domain knowledge informs 15 template motifs (8 exact, 7 parametric) capturing well-known quantum circuit primitives:

\begin{itemize}
    \item \textbf{Phase gadgets}: A T-like Z-spider connected via Hadamard edges to $k$ phaseless Z-spiders, implementing multi-qubit phase rotations.
    \item \textbf{CNOT pair}: The minimal Z--X spider pair via a simple edge, representing a controlled-NOT gate after spider fusion.
    \item \textbf{ZZ interaction}: A Z--Z--Z chain with an arbitrary-phase center, the core of QAOA and Ising simulation circuits.
    \item \textbf{Syndrome extraction}: A Z-spider fan-out to multiple X-spiders, ubiquitous in error-correcting codes.
    \item \textbf{Toffoli core}: A T--Z--T chain from the standard Toffoli decomposition into Clifford+T gates.
    \item \textbf{Cluster chain}: Three Z-spiders connected by Hadamard edges, the building block of measurement-based quantum computation.
\end{itemize}

Parametric variants use phase wildcards (\texttt{any}, \texttt{any\_nonzero}, \texttt{any\_nonclifford}) to match structural families rather than exact phase assignments.

\subsection{Bottom-Up Enumeration}

We enumerate all connected subgraphs of size 3--5 across the corpus at the \texttt{spider\_fused} level.
Starting from each non-boundary node, a recursive BFS expands subgraphs up to the size limit, capping at 500 candidates per host graph for tractability.

Each candidate is hashed using the Weisfeiler--Leman (WL) algorithm~\cite{Weisfeiler1968} with three iterations, incorporating vertex type, phase class, and edge type into node labels.
Candidates appearing in $\geq 2$ distinct algorithms are retained; hash collisions are resolved by VF2 isomorphism confirmation.

\subsection{Neighborhood Extraction}

We extract radius-2 neighborhoods around ``interesting'' nodes satisfying one of three criteria:
\begin{enumerate}
    \item \textbf{Non-Clifford}: Nodes with t\_like or arbitrary phase class.
    \item \textbf{High-degree}: Nodes with degree $> 4$.
    \item \textbf{Color boundary}: Nodes adjacent to a different spider color.
\end{enumerate}

These neighborhoods capture the local context of structurally significant features.
After WL deduplication and cross-algorithm frequency filtering, recurring neighborhoods become motif candidates.

\subsection{Deduplication and Library Optimization}

The three strategies yield a combined set of 697 unique motifs after deduplication.
A weighted greedy set-cover optimization selects the most informative subset for fingerprinting, scoring each candidate motif by:
\begin{equation}
    w_i = \Delta_{\text{cov}}(i) + \lambda_d \cdot \Delta_{\text{div}}(i) + \lambda_s \cdot s(i),
\end{equation}
where $\Delta_{\text{cov}}$ is the marginal vertex coverage, $\Delta_{\text{div}}$ is the algorithm diversity gain (fraction of newly covered algorithms), $s(i)$ is a size preference, and $\lambda_d = 0.3$, $\lambda_s = 0.1$ are tunable weights.

% ============================================================================
\section{Fingerprinting and Phylogenetic Analysis}
\label{sec:fingerprinting}

\subsection{Fingerprint Matrix Construction}

For a corpus of $M$ algorithms and a motif library of $N$ motifs, we construct an $M \times N$ fingerprint matrix $\mathbf{F}$ where $F_{ij}$ counts the number of occurrences of motif $j$ in the ZX diagram of algorithm $i$ at the \texttt{spider\_fused} simplification level.
Row-normalizing $\mathbf{F}$ to unit $\ell_1$ norm yields a frequency matrix $\hat{\mathbf{F}}$ representing the motif \emph{profile} of each algorithm.

\subsection{Hierarchical Clustering}

Hierarchical agglomerative clustering on the cosine distance $d(\hat{\mathbf{f}}_i, \hat{\mathbf{f}}_j) = 1 - \cos(\hat{\mathbf{f}}_i, \hat{\mathbf{f}}_j)$ produces a dendrogram (motif phylogeny) that recovers known algorithmic family structure.
PCA projection into two dimensions reveals that the first two principal components explain approximately 25\% of variance, with meaningful family-level clustering.

\subsection{Motif Universality Spectrum}

Analysis of motif prevalence across families reveals a universality spectrum:
\begin{itemize}
    \item \textbf{6 universal motifs} ($\geq$80\% family coverage): These represent fundamental ZX primitives shared across quantum computing, including the CNOT spider pair and basic phase gadgets.
    \item \textbf{30 common motifs} ($\geq$40\% coverage): Cross-domain structural elements.
    \item \textbf{655 family-specific motifs}: Characteristic signatures of particular algorithmic families.
    \item \textbf{6 unused motifs}: Template motifs not matched in the current corpus.
\end{itemize}

Decomposition analysis shows a mean motif coverage of 78.6\% across all algorithms, with error correction achieving the highest coverage (90.3\%) and distillation the lowest (41.7\%), suggesting that distillation circuits possess more unique structural elements not captured by the current library.

% ============================================================================
\section{Cross-Level Motif Evolution}
\label{sec:cross_level}

A distinctive feature of our analysis is tracking motif presence across all six simplification levels.
For each motif $m$ and algorithm $a$, we classify the motif's fate at each level $\ell$ as:
\begin{itemize}
    \item \textbf{Survives}: Exact match found at level $\ell$.
    \item \textbf{Transforms}: No exact match, but a structurally similar subgraph exists (cosine similarity $\geq 0.7$ on a 12-dimensional feature vector).
    \item \textbf{Vanishes}: Neither exact nor approximate match.
\end{itemize}

Of the full motif library, 10 motifs survive \texttt{full\_reduce}---we term these \emph{irreducible motifs}.
Irreducible motifs represent structural elements that cannot be further simplified by any ZX rewrite rule, making them fundamental building blocks of quantum computation in the ZX representation.

% ============================================================================
\section{Structure-Guided Ansatz Discovery}
\label{sec:ansatz_discovery}

\subsection{Irreducible Composition}

We exploit irreducible motifs for automated circuit design.
The key insight is that motifs surviving full ZX reduction encode computational primitives that are ``maximally compressed''---they resist simplification because they perform irreducible quantum operations.

Our discovery strategy composes pairs of irreducible motifs into candidate ansatz circuits by:
\begin{enumerate}
    \item Identifying all irreducible motifs from cross-level analysis.
    \item Generating pairwise sequential and parallel compositions using the ZX box algebra (Section~\ref{sec:composition}).
    \item Validating unitarity and computing novelty scores against the existing corpus.
    \item Evaluating promising candidates via VQE benchmarks.
\end{enumerate}

This process yielded six candidate ansätze (\texttt{irr\_pair01} through \texttt{irr\_pair06}), of which \texttt{irr\_pair11} emerged as the most promising after initial screening.

\subsection{The \texttt{irr\_pair11} Ansatz}
\label{sec:irr_pair11}

The \texttt{irr\_pair11} ansatz has three structural components, directly reflecting its motif origins:

\begin{enumerate}
    \item \textbf{Star hub}: Qubit 0 serves as a central entanglement hub, with CX gates fanning out to qubits $1, \ldots, k$ (where $k$ scales sublinearly with $n$). This originates from the graph-state star motif.
    \item \textbf{Phase gadgets}: T gates on every third qubit, conjugated by CX pairs, creating irreducible non-Clifford entanglement. This originates from the phase-gadget motif.
    \item \textbf{Chain tail}: Nearest-neighbor CX gates extending entanglement to remaining qubits. This provides linear-depth connectivity.
\end{enumerate}

The circuit scales linearly with qubit count and uses $O(n)$ two-qubit gates.
Under full ZX simplification, the diagram compresses by 19--26\%, confirming non-trivial but irreducible structure.

\subsection{ZX Box Composition}
\label{sec:composition}

Motif composition is formalized through a \emph{ZX box algebra}.
Each motif instance is wrapped as a \texttt{ZXBox} with explicit left (input) and right (output) boundaries.
Two composition operators are defined:

\textbf{Sequential composition} ($A ; B$): Identifies the output boundary of $A$ with the input boundary of $B$, fusing corresponding boundary spiders.

\textbf{Parallel composition} ($A \otimes B$): Tensor product placement, concatenating boundaries.

Simplification of composed boxes applies PyZX reduction rules while monitoring boundary preservation; if a higher-level simplification would destroy boundaries, it is rejected (raising a \texttt{BoundaryDestroyedError}).

\subsection{VQE Evaluation}

We benchmark \texttt{irr\_pair11} against two baselines with matched two-qubit gate budgets:
\begin{itemize}
    \item \textbf{CX-chain}: Linear nearest-neighbor CNOT chain.
    \item \textbf{HEA}: Hardware-efficient ansatz with brick-layer CZ gates.
\end{itemize}

Each configuration uses an RY/RZ--entangler--RY/RZ structure with $4n$ variational parameters.
We evaluate on three Hamiltonian models (Heisenberg, transverse-field Ising, XY) at two qubit counts, with 10 restarts per seed and 5 independent seeds per configuration.

\begin{table}[h]
\caption{VQE relative error (\%) for \texttt{irr\_pair11} versus baselines on 6-qubit systems. Lower is better. All differences are statistically significant ($p < 0.001$).}
\label{tab:vqe}
\begin{ruledtabular}
\begin{tabular}{lccc}
\textbf{Hamiltonian} & \textbf{irr\_pair11} & \textbf{CX-chain} & \textbf{HEA} \\
\midrule
Heisenberg & 23.9 & 25.0 & 24.8 \\
TFIM & 3.85 & 3.96 & 3.99 \\
XY & --- & --- & --- \\
\end{tabular}
\end{ruledtabular}
\end{table}

The \texttt{irr\_pair11} ansatz wins on all six benchmark configurations, with Cohen's $d$ up to 450, indicating very large effect sizes despite modest absolute improvements.
The practical significance lies in the systematic derivation: this ansatz was discovered automatically from structural analysis, not manual engineering.

\subsection{Ablation Study}

To confirm that each structural component is necessary, we evaluate ablated versions:
\begin{itemize}
    \item \textbf{No star hub}: Removes the central CX fan-out. Relative error increases by 27\%.
    \item \textbf{No phase gadgets}: Removes T-gate conjugation. Relative error increases by 52\%.
    \item \textbf{No chain tail}: Removes linear connectivity. No impact at 6 qubits (redundant at small scale), but expected to matter at larger qubit counts.
\end{itemize}

Phase gadgets contribute the most, consistent with the hypothesis that irreducible non-Clifford elements are critical for variational expressibility.

% ============================================================================
\section{Gap-Driven Motif Induction}
\label{sec:induction}

To bootstrap the motif library beyond initial discovery, we implement an iterative induction loop.
After decomposing each algorithm into non-overlapping motif placements via greedy set-cover, uncovered vertices are analyzed:

\begin{enumerate}
    \item \textbf{Signature clustering}: Each uncovered vertex is assigned a 1-hop structural signature encoding its type, phase, neighbor types, and degree.
    \item \textbf{Frequency ranking}: Signatures are ranked by frequency $\times$ algorithm diversity.
    \item \textbf{Extraction}: Radius-1 neighborhoods of high-ranking vertices are extracted as candidate motifs.
    \item \textbf{Validation}: Candidates are deduplicated (WL hash + VF2) and verified against the corpus ($\geq$5 occurrences across $\geq$2 algorithms).
\end{enumerate}

Iterative application (up to 3 rounds, stopping when coverage gain $< 2\%$) progressively fills gaps in the motif library.

% ============================================================================
\section{Implementation}
\label{sec:implementation}

\zxmotifs{} is implemented in Python and built on three core dependencies: Qiskit~\cite{Qiskit2024} for circuit definition, PyZX~\cite{Kissinger2020} for ZX-calculus operations, and NetworkX~\cite{Hagberg2008} for graph algorithms.
The system is organized as an installable package with a command-line interface supporting algorithm/motif listing, information queries, scaffolding for new contributions, and full-corpus validation.

All algorithms are registered via a decorator pattern:
\begin{verbatim}
@register_algorithm(
    name="grover",
    family="oracle",
    qubit_range=(3, 6))
def make_grover(n_qubits=3, **kwargs):
    ...
    return QuantumCircuit(...)
\end{verbatim}

Motifs are defined declaratively as JSON files validated against a JSON Schema, ensuring consistency and enabling version-controlled library evolution.

The full test suite comprises 200+ tests validating every registered algorithm, every motif pattern, and every pipeline stage.
Continuous integration tests on Python 3.10 and 3.12.

% ============================================================================
\section{Related Work}
\label{sec:related}

\textbf{ZX-calculus for optimization.}
PyZX~\cite{Kissinger2020} pioneered automated circuit optimization via ZX rewrite rules.
Our work complements PyZX by using ZX diagrams as an \emph{analysis} substrate rather than an optimization target.

\textbf{Quantum circuit patterns.}
Prior work on quantum circuit patterns has focused on gate-level template matching for peephole optimization~\cite{Nam2018} or equivalence checking~\cite{Amy2019}.
Our approach operates at the ZX-diagram level, where gate-level distinctions are abstracted away and deeper structural relationships become visible.

\textbf{Subgraph mining.}
Frequent subgraph mining is well-studied in bioinformatics and social networks~\cite{Yan2002}.
Our adaptation to ZX diagrams introduces domain-specific challenges: semantic node/edge labels, phase wildcards, and the multi-level simplification hierarchy.

\textbf{Ansatz design.}
Variational ansatz design has been approached through hardware-efficient strategies~\cite{Kandala2017}, ADAPT-VQE~\cite{Grimsley2019}, and reinforcement learning~\cite{Kuo2021}.
Our structure-guided approach is, to our knowledge, the first to derive ansätze from ZX-calculus motif composition.

% ============================================================================
\section{Discussion and Future Work}
\label{sec:discussion}

\textbf{Scalability.}
The current corpus is limited to $\leq$8 qubits for most algorithms due to VF2 computational cost.
Approximate matching with bounded edit distance partially addresses this, but scaling to industrially relevant circuit sizes will require more efficient matching algorithms or graph neural network surrogates.

\textbf{Motif completeness.}
Our 697-motif library achieves 78.6\% mean coverage, leaving room for improvement.
The iterative induction loop provides a principled path toward higher coverage, but fundamental questions about motif library completeness remain open.

\textbf{Noise awareness.}
The current analysis operates on ideal circuits.
Extending to noise-aware ZX diagrams---incorporating quantum channels as CP maps in the ZX framework---would enable motif-based analysis of error propagation and mitigation strategies.

\textbf{Automated discovery at scale.}
The success of \texttt{irr\_pair11} suggests that systematic motif composition could yield a rich space of novel ansätze.
Automated exploration of this space, guided by motif-level properties rather than gate-level search, is a promising direction.

% ============================================================================
\section{Conclusion}
\label{sec:conclusion}

We have presented \zxmotifs{}, a framework for quantum algorithm analysis through the lens of ZX-calculus structural motifs.
By converting 78 algorithms into semantically labeled ZX diagrams, discovering 697 recurring motifs via three complementary strategies, and constructing fingerprint vectors, we demonstrate that motif profiles capture meaningful algorithmic structure---recovering known family relationships and enabling cross-family structural comparison.

The identification of irreducible motifs and their composition into novel variational ansätze represents a new paradigm for structure-guided quantum circuit design.
The statistically significant VQE improvements achieved by \texttt{irr\_pair11}, derived entirely from structural analysis rather than manual engineering, validate this approach.

The \zxmotifs{} framework, including the full algorithm corpus, motif library, and analysis pipeline, is available as open-source software at \url{https://github.com/splch/zx-motifs}.

\begin{acknowledgments}
We thank the developers of PyZX, Qiskit, and NetworkX for the foundational tools that make this work possible.
\end{acknowledgments}

\bibliography{references}

\end{document}
